%%
%% fall-lecture08.tex
%% 
%% Made by Alex Nelson <pqnelson@gmail.com>
%% Login   <alex@lisp>
%% 
%% Started on  2026-10-15T08:56:47-0700
%% Last update 2026-10-15T08:56:47-0700
%% 

\lecture{}

(Emma Lenz is lecturing today, the professor was out of town.)

\begin{recall}
We have derived the equations of motion for various quantities. Let us
review them, so they are all in one handy location.
\begin{itemize}
\item Momentum $\displaystyle{\rho\frac{D u_{i}}{D t} = -\partial_{i}p+\partial_{j}\tau_{ij}+\gamma_{i}}$
\item Kinetic energy $\displaystyle{\rho\frac{D(\frac{1}{2}u_{i}u_{i})}{D t} = -u_{i}\partial_{i}p+u_{i}\partial_{j}\tau_{ij}+u_{i}\gamma_{i}}$
\item Total energy (which is basically the first Law of Thermodynamics)
\begin{equation*}
\rho\frac{D\varepsilon}{Dt}=-\partial_{i}(pu_{i})+\partial_{j}(\tau_{ij}u_{i})+\gamma_{i}u_{i}-\partial_{i}q_{i}+Q\rho
\end{equation*}
where $\varepsilon=e+\frac{1}{2}u_{i}u_{i}$, and $e$ is thhe internal energy
\item Energy dissipation (obtained by subtracting the kinetic energy
  from the total energy):
\begin{align*}
\rho\frac{De}{Dt} &= \rho\frac{D}{Dt}(\varepsilon-\frac{1}{2}u_{i}u_{i})\\
&=-p\partial_{i}u_{i}+\tau_{ij}\partial_{j}u_{i}-\partial_{i}q_{i}+\rho Q
\end{align*}
\item Enthalpy $h=e+(p/\rho)$:
\begin{align*}
\rho\frac{Dh}{Dt} &= \rho\frac{D e}{Dt}+\rho\frac{D(p/\rho)}{Dt}\\
&=\tau_{ij}\partial_{j}u_{i}-\partial_{i}q_{i}+\rho Q+\partial_{t}p+u_{i}\partial_{i}p
\end{align*}
\item Total Enthalpy (a.k.a., ``Stagnation Enthalpy'') $h_{0}=h+\frac{1}{2}u_{i}u_{i}$
\begin{align*}
\rho\frac{Dh_{0}}{Dt} &= \rho\frac{D h}{Dt}+\rho\frac{D(\frac{1}{2}u_{i}u_{i})}{Dt}\\
&=\partial_{t}p+\partial_{j}(\tau_{ij}u_{i})+u_{i}\gamma_{i}-\partial_{i}q_{i}+\rho Q
\end{align*}
\end{itemize}
A lot of these are just different ways to express the conservation of energy.
\end{recall}

\subsection{Fluid statics}

\begin{definition}
We call the situation where the fluid's velocity vector field vanishes
everywhere $\vec{u}=0$ \define{Fluid Statics}.
\end{definition}

\begin{node}
Observe in fluid statics $\tau_{ij}=0$ since nothing is moving, so
nothing can generate shear.
\end{node}

\begin{node}
The conservation of linear momentum becomes
\begin{subequations}
\begin{equation}
\underbrace{\rho\frac{D u_{i}}{D t}}_{=0} = -\partial_{i}p+\underbrace{\partial_{j}\tau_{ij}}_{=0}+\gamma_{i},
\end{equation}
which can be re-arranged to
\begin{equation}
\partial_{i}p=\gamma_{i}.
\end{equation}
\end{subequations}
In other words, the pressure gradient $\partial_{i}p$ balances body
forces $\gamma_{i}$.
\end{node}

\begin{node}
When the external force is the Earth's gravitational field,
$\vec{\gamma}=\rho\vec{g}=-\rho g\widehat{k}$, then we have
\begin{equation}
\partial_{z}p=-\rho(z,p) g
\end{equation}
and the other two equations vanish $\partial_{x}p=\partial_{y}p=0$.
\end{node}

\begin{node}[Ideal gases]
For ideal gases,
\begin{equation}
\rho=\frac{p}{RT}
\end{equation}
which gives us
\begin{equation}
\partial_{z}p=-\frac{pg}{RT},
\end{equation}
or
\begin{equation}
\frac{\D p}{p}=\frac{-g\,\D z}{RT(z)}.
\end{equation}
We only need to know the temperature $T(z)$ to infer anything
further. For Earth's atmosphere, $T$ linearly decreases in the
troposphere until it gets to the tropopause, and then is constant (and
then increases, and then\dots).
\end{node}

\begin{node}
(This was not discussed in class.) If we take $T(z)=T_{0}-T'z$, then
(assuming $T_{0}>T'z$):
\begin{equation}
\int^{z}_{0}\frac{-g}{R(T_{0}-T'z')}\D z'=\frac{-g\ln(1 - (T'/T_{0})z)}{T'R},
\end{equation}
which gives us
\begin{equation}
\ln(p(z)/p(0))=\frac{-g\ln\left(1 - \frac{T'z}{T_{0}}\right)}{T'R},
\end{equation}
or equivalently,
\begin{equation}
p(z) = p_{0}\left(1 - \frac{T'z}{T_{0}}\right)^{g/(T'R)}\sim p_{0}(1-cz)^{b}
\end{equation}
where $c$, $b$ are positive real constants.
\end{node}

\subsection{Uniform acceleration}

\begin{example}[Linear acceleration]
If we have a box half full of liquid (or any fluid, really), then we
can draw the free body diagram for a fluid parcel at time $t=0$
(before a uniform linear acceleration is applied):
\begin{center}
\includegraphics{img/img.10}
\end{center}
Now, if the container is uniformly accelerating ``forward'' with
acceleration $a$, then at a later time we would expect the following
free body diagram:
\begin{center}
\includegraphics{img/img.11}
\end{center}
where $\gamma_{\text{eff}}$ is the effective force acting on the fluid parcel.
This must be balanced by the pressure gradient, which has changed accordingly.
\end{example}

\begin{example}[Uniform solid body rotation]
Now, consider the situation where the container is a
cylinder. Initially, at time $t=0$, it is at rest and the cross
section looks something like:
\begin{center}
  \includegraphics{img/img.12}
\end{center}
Then we can consider the situation where it rotates with constant
angular velocity $\omega$ about the $z$-axis. This will have the free
body diagram:
\begin{center}
  \includegraphics{img/img.13}
\end{center}
Again, the pressure gradient balances the effective force
$\vec{\gamma}_{\text{eff}}$. In cylindrical coordinates, the forces
acting on the fluid parcel are:
\begin{equation}
\partial_{z}p\widehat{z}+\partial_{r}p\widehat{r}=-\rho g\widehat{z}+\rho\omega^{2}r\widehat{r}
\end{equation}
This example has a particularly fascinating application building
telescopes using reflective fluids [instead of glass] rotating around
at a constant angular velocity $\omega$ (namely, at different angular speeds
$\omega$, the focal length of the telescope changes).
\end{example}

\subsection{Euler and Bernoulli equations}

\begin{assume}
The assumptions for the Euler equations:
\begin{enumerate}
\item No shear stress $\tau_{ij}=0$
\item No heat transfer $q_{i}=0$
\item The fluid is an ideal inviscid fluid
\end{enumerate}
\end{assume}

\begin{node}
The equations of motion simplify considerably, and we refer to them as
\define{Euler's Equations}.\footnote{Observe the continuity equation
and the conservation of linear momentum gives us 4 equations in 5
unknowns. Usually an extra equation is given relating pressure to
density, justified by Thermodynamics.} The continuity equation
becomes
\begin{equation}
\frac{D\rho}{Dt}=-\rho\partial_{i}u_{i}.
\end{equation}
The linear momentum equation becomes
\begin{subequations}
\begin{equation}
\rho\frac{D u_{i}}{D t}=-\partial_{i}p+\gamma_{i}
\end{equation}
or equivalently,
\begin{equation}
\partial_{t}(\rho u_{i})=-\partial_{j}(\rho u_{i}u_{j})-\partial_{i}p-\gamma_{i}.
\end{equation}
\end{subequations}
The conservation of energy becomes
\begin{equation}
\rho\frac{D\varepsilon}{D t}=-\partial_{i}(pu_{i})+\gamma_{i}u_{i}+\rho Q.
\end{equation}
The conservation of total enthalpy becomes:
\begin{equation}
\rho\frac{D h_{0}}{D t}=-\partial_{t}p+\rho Q+\gamma_{i}u_{i}.
\end{equation}
We can use these to deduce Bernoulli's principle.
\end{node}

\subsection*{Deriving Bernoulli's principle from Euler's equations}

\begin{assume}
We will assume:
\begin{enumerate}
\item No shear stress $\tau_{ij}=0$
\item No heat transfer $q_{i}=0$
\item The fluid is an ideal inviscid fluid
\item Steady flow: $\partial_{t}\mbox{(anything)}=0$
\item No heat source $Q=0$
\end{enumerate}
Observe the first three assumptions merely allow us to use Euler's
equations.

We will also define a force potential $V$ [scalar function] such that
the body force satisfies:
\begin{equation}
\vec{\gamma}=-\rho\vec{\nabla}V.
\end{equation}
\end{assume}

\begin{node}[Bernoulli's principle for (possibly compressible) flows]
From the total Enthalpy equation
\begin{subequations}
  \begin{align}
\rho\frac{Dh_{0}}{Dt}
&=-\partial_{t}p+\rho Q+\gamma_{i}u_{i}\\
&=0+\rho Q +\gamma_{i}u_{i}\quad\mbox{by steady flow}\\
&=0+\gamma_{i}u_{i}\quad\mbox{by no heat source}\\
&=(-\rho\partial_{i}V)u_{i}\\
&=-\rho\left(\frac{DV}{Dt}-\partial_{t}V\right)
  \end{align}
\end{subequations}
where the last line takes advantage of steady flow. Then we rearrange
terms to find:
\begin{subequations}
  \begin{align}
\rho\left(\frac{Dh_{0}}{Dt}+\frac{DV}{Dt}\right)
&=\rho\frac{D}{Dt}(h_{0}+V)\\
&=\rho\frac{D}{Dt}(h_{0}+\frac{1}{2}u_{i}u_{i}+V)\quad\mbox{by steady flow}\\
&=0.
  \end{align}
\end{subequations}
Since $\rho$ is nonzero, we can divide it out, and conclude
\begin{equation}
\frac{D}{Dt}(h_{0}+\frac{1}{2}u_{i}u_{i}+V)=0.
\end{equation}
What does it mean for a material derivative of a function to be zero?
Along the path line [trajectory], the quantity is constant.
\end{node}

\begin{definition}\label{defn:stream-line}
\define{Stream Lines} are curves which are tangent to the velocity vector.
They satisfy
\begin{equation*}
\frac{\D x}{u}=\frac{\D y}{v}=\frac{\D z}{w}
\end{equation*}
where $u=\widehat{x}\cdot\vec{u}$, $v=\widehat{y}\cdot\vec{u}$, $w=\widehat{z}\cdot\vec{u}$.
\end{definition}

\begin{definition}
\define{Path Lines} are physical trajectories of ``fluid parcels''.
\end{definition}

\begin{fact}
In a steady flow, stream lines coincide with path lines.
\end{fact}

\begin{node}[Bernoulli's principle for incompressible flows]
Furthermore, for incompressible flows, we have
\begin{equation}
\frac{\D\rho}{Dt}=0\quad\mbox{implies}\quad\vec{\nabla}\cdot\vec{u}=0.
\end{equation}
If we add this to our running list of assumptions, then we infer
\begin{equation}
\frac{De}{Dt}=0,
\end{equation}
so the internal energy $e$ is constant along stream lines. So we find
\begin{equation}
h + \frac{1}{2}u_{i}u_{i}+V=\left(e +
\frac{p}{\rho}\right)+\frac{1}{2}u_{i}u_{i}+V=\mbox{constant along streamlines},
\end{equation}
which implies (since $e$ is constant along streamlines)
\begin{equation}
\frac{p}{\rho}+\frac{1}{2}u_{i}u_{i}+V=\mbox{constant along streamlines}.
\end{equation}
Rearranging terms (and since $\rho$ is constant along stream lines, by
continuity equation)
\begin{equation}
p+\frac{1}{2}\rho u_{i}u_{i}+\rho V=\mbox{constant along streamlines}.
\end{equation}
When gravity is the only external force, $\vec{\gamma}=\rho\vec{g}$
and $V=gz$ gives us
\begin{equation}
p+\frac{1}{2}\rho u_{i}u_{i}+\rho gz=\mbox{constant along streamlines}.
\end{equation}
This is the familiar form of Bernoulli's equations for incompressible flows.
\end{node}

\begin{remark}
Importantly, we have demonstrated the ``basic process'' for deducing
Bernoulli's principle from the equations of motion. We can
\emph{generalize} Bernoulli's principle to other settings by
performing the same calculations with weaker assumptions.
\end{remark}